% ============================================================================
% CHAPTER 5: TRANSFORMERS AND MODERN ARCHITECTURES
% The single most important chapter for NLP interviews.
% ============================================================================

\chapter{Transformers and Modern Architectures}
\label{chap:transformers}

\begin{tldr}
The Transformer architecture (Vaswani et al., 2017) revolutionized NLP by replacing
recurrence with \emph{self-attention}, enabling full parallelism over the input sequence
and providing direct paths between any two positions.  This chapter derives
self-attention from first principles, builds up multi-head attention, positional
encoding, and the complete encoder--decoder architecture, then surveys the three
dominant variants---encoder-only (BERT), decoder-only (GPT), and encoder--decoder
(T5)---and modern efficiency improvements (Flash Attention, GQA, sliding window).
Transformers are the \textbf{\#1 most-tested topic in NLP interviews}: every
FAANG-level ML interview will probe your understanding of the math, the design
choices, and the production implications.
\end{tldr}

\bigskip

% ============================================================================
% SECTION 1: SELF-ATTENTION
% ============================================================================
\section{Self-Attention from First Principles}
\label{sec:self-attention}

\subsection{Motivation}

Recurrent models process tokens sequentially: to connect the first and last
tokens in a length-$n$ sequence, information must pass through $O(n)$
intermediate hidden states.  This sequential bottleneck causes two problems:

\begin{enumerate}
    \item \textbf{Training is slow.}  Backpropagation through time cannot be
          parallelized across positions; each step depends on the previous one.
    \item \textbf{Long-range dependencies are hard to learn.}  Even with LSTM
          gating, gradient signals attenuate over many steps, making it
          difficult to capture dependencies spanning hundreds of tokens.
\end{enumerate}

Self-attention solves both problems by computing \emph{direct} pairwise
interactions between all positions in a single operation, reducing the path
length between any two tokens to $O(1)$ and enabling full parallelism.

\subsection{Deriving Self-Attention Step by Step}

Suppose we have a sequence of $n$ token representations, each of dimension $d$,
stacked into a matrix:
\[
    \mX \in \R^{n \times d}
\]
where row $\mX_i$ is the representation of the $i$-th token.

\paragraph{Step 1: Project into Query, Key, and Value spaces.}
We learn three separate linear projections:
\[
    \mQ = \mX \mW_Q, \qquad
    \mK = \mX \mW_K, \qquad
    \mV = \mX \mW_V
\]
where $\mW_Q, \mW_K \in \R^{d \times d_k}$ and $\mW_V \in \R^{d \times d_v}$.

\paragraph{Step 2: Compute attention scores.}
For every pair of positions $(i,j)$, compute a compatibility score between
query $i$ and key $j$ via a scaled dot product:
\[
    e_{ij} = \frac{\vq_i^\top \vk_j}{\sqrt{d_k}}
\]

\paragraph{Step 3: Normalize into attention weights.}
Apply $\softmax$ row-wise so that each query's weights sum to one:
\[
    \alpha_{ij} = \frac{\exp(e_{ij})}{\sum_{l=1}^{n} \exp(e_{il})}
\]

\paragraph{Step 4: Compute weighted sum of values.}
The output for position $i$ is the attention-weighted combination of all values:
\[
    \text{output}_i = \sum_{j=1}^{n} \alpha_{ij}\, \mathbf{v}_j
\]

In full matrix form, this gives the celebrated attention equation:

\begin{mathresult}{Scaled Dot-Product Attention}
\[
    \text{Attention}(\mQ, \mK, \mV)
    = \softmax\!\left(\frac{\mQ \mK^\top}{\sqrt{d_k}}\right) \mV
\]
where $\mQ \in \R^{n \times d_k}$, $\mK \in \R^{n \times d_k}$,
$\mV \in \R^{n \times d_v}$, and $\softmax$ is applied row-wise.
\end{mathresult}

\subsection{Why Separate Q, K, V Projections?}

A common interview follow-up is: \emph{why not just use $\mX$ directly for the
dot products?}

Without separate projections, the attention score between positions $i$ and $j$
would be $\vx_i^\top \vx_j$---a simple input similarity measure.  This is
degenerate: for representations of approximately equal norm, every token
would attend most strongly to itself (Cauchy--Schwarz gives
$\vx_i^\top \vx_j \leq \norm{\vx_i}\,\norm{\vx_j} \approx \norm{\vx_i}^2
= \vx_i^\top \vx_i$, with equality only when $\vx_j$ is parallel to
$\vx_i$).

Separate projections decouple three distinct \emph{roles}:
\begin{itemize}
    \item $\mQ$: \textbf{``What am I looking for?''} --- the query's
          information need.
    \item $\mK$: \textbf{``What do I advertise?''} --- what each position
          offers as a potential match.
    \item $\mV$: \textbf{``What do I provide if selected?''} --- the content
          actually transmitted.
\end{itemize}

This separation allows the model to attend based on one aspect of the
representation (e.g., syntactic role via Q/K) while retrieving a different
aspect (e.g., semantic content via V).

\begin{keyinsight}[Attention as Soft Dictionary Lookup]
Self-attention is a \emph{soft dictionary lookup}.  The queries search over
keys to produce a probability distribution, then retrieve a weighted mixture
of the associated values.  Unlike a hard lookup (which returns a single
entry), this soft version is differentiable and can blend information from
multiple positions simultaneously.  This analogy is directly useful for
understanding the KV-cache during inference, where keys and values from
previously generated tokens are stored for reuse.
\end{keyinsight}

\subsection{Why Scale by $\sqrt{d_k}$?}
\label{sec:sqrt-scaling}

This is one of the most commonly asked interview questions.  The answer
requires a simple variance analysis.

Assume each component of $\vq$ and $\vk$ is drawn i.i.d.\ with mean 0 and
variance 1.  The dot product is:
\[
    \vq^\top \vk = \sum_{i=1}^{d_k} q_i \, k_i
\]

Since $q_i$ and $k_i$ are independent, $\E[q_i k_i] = 0$ and
$\Var(q_i k_i) = \E[q_i^2]\,\E[k_i^2] = 1$.  By independence of the summands:

\begin{mathresult}{Variance of the Dot Product}
\[
    \Var(\vq^\top \vk) = \sum_{i=1}^{d_k} \Var(q_i\, k_i) = d_k
\]
\end{mathresult}

So the standard deviation of the raw dot product grows as $\sqrt{d_k}$.  For
typical $d_k = 64$ or $128$, the dot products range over $\pm 16$ or more.
The $\softmax$ function saturates in this regime: its outputs approach a one-hot
vector, gradients vanish, and the model cannot learn.

Dividing by $\sqrt{d_k}$ normalizes the variance to 1 regardless of the
dimension, keeping the $\softmax$ in its ``useful'' regime where gradients
flow meaningfully.

\begin{warningbox}
Candidates frequently state ``we divide by $\sqrt{d_k}$ to normalize'' without
being able to derive \emph{why} $\sqrt{d_k}$.  Be prepared to walk through the
variance calculation above---interviewers will push on it.
\end{warningbox}

\subsection{Computational Complexity}

The matrix product $\mQ \mK^\top$ produces an $n \times n$ matrix, where $n$
is the sequence length.  This yields:

\begin{itemize}
    \item \textbf{Compute:} $O(n^2 d)$ for the QK and attention-value products.
    \item \textbf{Memory:} $O(n^2)$ for the attention matrix (plus $O(nd)$
          for Q, K, V).
\end{itemize}

For $n = 4096$ and 32 heads, the attention matrix alone is
$32 \times 4096^2 \times 2\text{ bytes} \approx 1$\,GB in \texttt{float16}.
This quadratic scaling is the fundamental bottleneck driving all efficiency
research (Flash Attention, sparse attention, linear attention, SSMs).

\subsection{TikZ: Scaled Dot-Product Attention}

\begin{figure}[H]
\centering
\begin{tikzpicture}[
    node distance=1.0cm and 1.6cm,
    box/.style={rectangle, draw=deepblue, fill=lightblue, minimum width=2.0cm,
                minimum height=0.7cm, font=\small\sffamily, rounded corners=2pt,
                thick},
    opbox/.style={circle, draw=deepblue, fill=deepblue!12, minimum size=0.7cm,
                  font=\small\sffamily, thick},
    arr/.style={-{Stealth[length=5pt]}, thick, deepblue!80},
    lbl/.style={font=\scriptsize\sffamily, text=deepblue!80}
]

% Input matrices
\node[box] (Q) {$\mQ$};
\node[box, right=1.8cm of Q] (K) {$\mK$};
\node[box, right=1.8cm of K] (V) {$\mV$};

% MatMul QK^T
\node[opbox, below=1.0cm of $(Q)!0.5!(K)$] (mm1) {$\times$};
\node[lbl, right=2pt of mm1] {$\mQ\mK^\top$};

% Scale
\node[box, below=0.8cm of mm1] (scale) {Scale $\div\!\sqrt{d_k}$};

% Mask (optional)
\node[box, below=0.8cm of scale, dashed, draw=gray] (mask) {Mask (opt.)};

% Softmax
\node[box, below=0.8cm of mask] (sm) {Softmax};

% MatMul with V
\node[opbox, below=0.8cm of $(sm)!0.5!(sm-|V)$] (mm2) {$\times$};

% Output
\node[box, below=0.8cm of mm2, fill=deepblue!15] (out) {Output};

% Arrows
\draw[arr] (Q) |- (mm1);
\draw[arr] (K) |- (mm1);
\draw[arr] (mm1) -- (scale);
\draw[arr] (scale) -- (mask);
\draw[arr] (mask) -- (sm);
\draw[arr] (sm) -- (mm2);
\draw[arr] (V) |- (mm2);
\draw[arr] (mm2) -- (out);

% Transpose label
\node[lbl, above left=0.1cm and -0.2cm of mm1] {$^\top$};

\end{tikzpicture}
\caption{Scaled dot-product attention.  Queries and keys interact via a scaled
dot product, optionally masked (for causal decoding), then normalized by
$\softmax$ to produce attention weights applied to values.}
\label{fig:sdp-attention}
\end{figure}


% ============================================================================
% SECTION 2: MULTI-HEAD ATTENTION
% ============================================================================
\section{Multi-Head Attention}
\label{sec:multi-head}

A single attention head captures one ``type'' of relationship (e.g., syntactic
adjacency).  \textbf{Multi-head attention} runs $h$ independent attention heads
in parallel, each with its own learned projections, then concatenates and
linearly combines the results.

\begin{mathresult}{Multi-Head Attention}
\begin{align*}
    \text{MultiHead}(\mQ, \mK, \mV) &= \text{Concat}(\text{head}_1, \ldots, \text{head}_h)\,\mW_O \\[4pt]
    \text{head}_i &= \text{Attention}(\mQ\mW_Q^{(i)},\; \mK\mW_K^{(i)},\; \mV\mW_V^{(i)})
\end{align*}
where $\mW_Q^{(i)}, \mW_K^{(i)} \in \R^{d_{\text{model}} \times d_k}$,\;
$\mW_V^{(i)} \in \R^{d_{\text{model}} \times d_v}$,\;
$\mW_O \in \R^{hd_v \times d_{\text{model}}}$, and
$d_k = d_v = d_{\text{model}} / h$.
\end{mathresult}

\paragraph{Why multiple heads?}
\begin{enumerate}
    \item \textbf{Diverse attention patterns.}  Different heads learn to focus
          on different relationship types---empirical studies show heads
          specializing in syntactic heads, positional neighbors, rare tokens,
          and semantic roles.
    \item \textbf{Subspace richness.}  Each head operates in its own
          $d_k$-dimensional subspace, allowing the model to attend based on
          multiple criteria simultaneously.
    \item \textbf{Ensemble effect.}  Multiple heads act as an ensemble within
          a single layer, increasing robustness.
\end{enumerate}

\paragraph{What happens with a single head?}
With one head at full dimension $d_k = d_{\text{model}}$, the model can
represent only \emph{one} attention pattern per layer.  Empirically,
single-head attention performs significantly worse, especially on tasks
requiring both syntactic and semantic reasoning.  The output projection
$\mW_O$ is also critical: it allows cross-head information mixing that a
single head cannot provide.

\paragraph{Computational cost.}
Multi-head attention has the \emph{same} total cost as single-head attention
at full dimension.  Each head uses $d_k = d_{\text{model}}/h$, so the total
parameter count and FLOPs remain $O(d_{\text{model}}^2)$ per layer---the
parallelism over heads is ``free.''

\begin{figure}[H]
\centering
\begin{tikzpicture}[
    node distance=0.8cm and 0.7cm,
    box/.style={rectangle, draw=deepblue, fill=lightblue, minimum width=1.4cm,
                minimum height=0.6cm, font=\scriptsize\sffamily, rounded corners=2pt,
                thick},
    headbox/.style={rectangle, draw=deepblue!70, fill=blue!6,
                    minimum width=1.3cm, minimum height=0.55cm,
                    font=\scriptsize\sffamily, rounded corners=2pt},
    arr/.style={-{Stealth[length=4pt]}, thick, deepblue!70},
    lbl/.style={font=\tiny\sffamily, text=deepblue!70}
]

% Input
\node[box, minimum width=5.2cm] (input) {Input: $\mX \in \R^{n \times d}$};

% Linear projections
\node[box, below=0.8cm of input, xshift=-2.0cm] (lq) {$\mW_Q^{(i)}$};
\node[box, below=0.8cm of input] (lk) {$\mW_K^{(i)}$};
\node[box, below=0.8cm of input, xshift=2.0cm] (lv) {$\mW_V^{(i)}$};

% Heads
\node[headbox, below=0.7cm of lq, xshift=-0.8cm] (h1) {Head 1};
\node[headbox, below=0.7cm of lk] (h2) {Head 2};
\node[lbl, right=0.15cm of h2] {$\cdots$};
\node[headbox, below=0.7cm of lv, xshift=0.8cm] (hh) {Head $h$};

% Concat
\node[box, below=0.7cm of h2, minimum width=5.2cm] (concat) {Concat};

% Output projection
\node[box, below=0.6cm of concat, minimum width=5.2cm, fill=deepblue!15] (wo) {Linear $\mW_O$};

% Output
\node[box, below=0.6cm of wo, minimum width=5.2cm, fill=deepblue!20] (out)
     {Output: $\R^{n \times d}$};

% Arrows
\draw[arr] (input) -- (lq);
\draw[arr] (input) -- (lk);
\draw[arr] (input) -- (lv);

\draw[arr] (lq) -- (h1);
\draw[arr] (lk) -- (h2);
\draw[arr] (lv) -- (hh);

\draw[arr] (h1) -- (concat);
\draw[arr] (h2) -- (concat);
\draw[arr] (hh) -- (concat);

\draw[arr] (concat) -- (wo);
\draw[arr] (wo) -- (out);

% Annotation
\node[lbl, left=0.3cm of h1] {$d_k = d/h$};

\end{tikzpicture}
\caption{Multi-head attention.  Input is projected $h$ times into lower-dimensional
subspaces, attention is computed independently per head, and results are
concatenated and linearly projected back to $d_{\text{model}}$ dimensions.}
\label{fig:multi-head}
\end{figure}


% ============================================================================
% SECTION 3: FEED-FORWARD NETWORK
% ============================================================================
\section{Position-Wise Feed-Forward Network}
\label{sec:ffn}

Each transformer layer contains a \emph{position-wise} feed-forward network
(FFN) applied identically and independently to every position.  Unlike
attention, the FFN involves \emph{no} cross-position interaction---it
transforms each token's representation in isolation.

\subsection{Standard FFN}

The original Transformer uses a two-layer MLP with a ReLU activation:
\[
    \text{FFN}(\vx) = \relu(\vx \mW_1 + \vb_1)\,\mW_2 + \vb_2
\]
where $\mW_1 \in \R^{d \times d_{ff}}$, $\mW_2 \in \R^{d_{ff} \times d}$,
and $d_{ff} = 4d$ is the standard expansion ratio.  Modern models often use
$\gelu$ instead of $\relu$ (used in BERT, GPT-2).

\subsection{Gated Linear Units: SwiGLU and GeGLU}

Most modern LLMs (LLaMA, Mistral, Gemma, PaLM) replace the standard FFN with
a \textbf{gated linear unit} variant:
\[
    \text{SwiGLU}(\vx) = \bigl(\text{Swish}(\vx \mW_1) \odot (\vx \mW_3)\bigr)\,\mW_2
\]
where $\text{Swish}(x) = x \cdot \sigmoid(x)$, $\odot$ is element-wise
multiplication, and $\mW_3 \in \R^{d \times d_{ff}}$ is the gate projection.
The gating mechanism allows the network to selectively pass or suppress
features, empirically improving performance for a given parameter count.

\paragraph{Why the expansion then compression?}
The FFN expands to $d_{ff} = 4d$ (or $\tfrac{8}{3}d$ in gated variants to
match parameter counts), then compresses back to $d$.  The high-dimensional
intermediate representation allows the network to project inputs into a space
where complex nonlinear transformations become easier, then distill the result
back to the model dimension.

\begin{keyinsight}[FFN as Key-Value Memory]
Geva et al.\ (2021) showed that FFN layers can be interpreted as
\emph{key-value memories}: the rows of $\mW_1$ are ``keys'' (pattern
detectors), the activation determines which patterns are triggered, and the
corresponding columns of $\mW_2$ are ``values'' (stored information retrieved
when the pattern matches).  This explains why scaling the FFN width improves
factual knowledge capacity and why knowledge editing techniques target FFN
layers specifically.
\end{keyinsight}

\paragraph{Attention vs.\ FFN: complementary roles.}
\begin{itemize}
    \item \textbf{Attention} = \emph{token mixing} (cross-position
          interaction, determines which positions to combine).
    \item \textbf{FFN} = \emph{feature mixing} (per-position transformation,
          applies nonlinear computations to each token's representation).
\end{itemize}

A transformer layer interleaves both: first attention combines information
across positions, then the FFN transforms each resulting representation.


% ============================================================================
% SECTION 4: POSITIONAL ENCODING
% ============================================================================
\section{Positional Encoding}
\label{sec:positional-encoding}

\subsection{Why Positional Information Is Needed}

Self-attention is \textbf{permutation-equivariant}: if we reorder the input
tokens, the output is reordered in exactly the same way, with the same values.
Formally, for any permutation $\pi$:
\[
    \text{Attention}(\mX_\pi) = \text{Attention}(\mX)_\pi
\]
This means the model has \emph{no innate notion of position}---``The cat sat
on the mat'' and ``mat the on sat cat The'' would produce identical attention
patterns.  Positional encoding injects position information to break this
symmetry.

\subsection{Sinusoidal Positional Encoding (Original Transformer)}

Vaswani et al.\ (2017) proposed a deterministic, parameter-free encoding using
sinusoids of varying frequencies:

\begin{mathresult}{Sinusoidal Positional Encoding}
\begin{align*}
    \text{PE}(pos, 2i) &= \sin\!\left(\frac{pos}{10000^{2i/d}}\right) \\[4pt]
    \text{PE}(pos, 2i+1) &= \cos\!\left(\frac{pos}{10000^{2i/d}}\right)
\end{align*}
where $pos$ is the position index and $i$ is the dimension index
($0 \leq i < d/2$).
\end{mathresult}

The key property: for any fixed offset $k$, $\text{PE}(pos + k)$ can be
expressed as a linear function of $\text{PE}(pos)$, enabling the model to
learn relative position attention.  The wavelengths form a geometric
progression from $2\pi$ to $10000 \cdot 2\pi$, covering fine-grained to
coarse-grained positional information.

\subsection{Learned Absolute Position Embeddings}

BERT, GPT-2, and many early models simply learn a position embedding table
$\mW_P \in \R^{L_{\max} \times d}$ and add $\mW_P[pos]$ to each token
embedding.

\textbf{Pros:} Simple, effective, can capture complex position-dependent patterns.\\
\textbf{Cons:} Fixed maximum length $L_{\max}$---cannot generalize to longer
sequences at inference time (no length extrapolation).

\subsection{Relative Position Encoding (Shaw et al., 2018)}

Instead of absolute positions, encode the \emph{relative distance} between
positions directly in the attention computation:
\[
    e_{ij} = \frac{\vq_i^\top(\vk_j + \va_{ij}^K)}{\sqrt{d_k}}
\]
where $\va_{ij}^K$ is a learned embedding for relative position $(i - j)$,
clipped to a maximum distance.  This naturally captures that ``two tokens apart''
matters more than ``at positions 57 and 59 specifically.''

\subsection{RoPE: Rotary Position Embedding}
\label{sec:rope}

Rotary Position Embedding (Su et al., 2021) encodes position by
\emph{rotating} query and key vectors in 2D subspaces:

\begin{mathresult}{Rotary Position Embedding (RoPE)}
For position $m$ and dimension pair $(2i, 2i+1)$:
\[
    \text{RoPE}(\vx, m) = \begin{pmatrix} x_{2i} \cos m\theta_i - x_{2i+1} \sin m\theta_i \\
                                           x_{2i} \sin m\theta_i + x_{2i+1} \cos m\theta_i \end{pmatrix}
\]
where $\theta_i = 10000^{-2i/d}$.  The inner product of rotated queries and
keys depends only on relative position:
\[
    \text{RoPE}(\vq, m)^\top \text{RoPE}(\vk, n) = f(\vq, \vk, m - n)
\]
\end{mathresult}

\textbf{Why RoPE is widely adopted} (LLaMA, Mistral, GPT-NeoX, Qwen, DeepSeek):
\begin{enumerate}
    \item Encodes \emph{relative} position through an \emph{absolute}
          mechanism---no separate relative position lookup.
    \item Decays naturally with distance: inner products decrease for faraway
          positions, providing an implicit recency bias.
    \item Supports length extrapolation through frequency scaling (NTK-aware
          scaling, YaRN, Dynamic NTK).
    \item No additional parameters beyond the base frequencies.
\end{enumerate}

\subsection{ALiBi: Attention with Linear Biases}

ALiBi (Press et al., 2022) takes an even simpler approach: subtract a linear
penalty from attention scores based on distance, with a different slope per head:
\[
    \text{score}_{ij} = \vq_i^\top \vk_j - m \cdot |i - j|
\]
where $m$ is a head-specific slope (set geometrically, not learned).  ALiBi
requires no position embeddings at all and achieves strong length extrapolation.

\subsection{Comparison of Positional Encoding Methods}

\begin{table}[H]
\centering
\small
\begin{tabularx}{\textwidth}{@{}l c c c X@{}}
\toprule
\textbf{Method} & \textbf{Params} & \textbf{Relative?} & \textbf{Extrapolation} & \textbf{Used In} \\
\midrule
Sinusoidal         & None    & Indirect & Limited  & Original Transformer \\
Learned Absolute   & $O(Ld)$ & No       & None     & BERT, GPT-2 \\
Relative (Shaw)    & $O(Kd)$ & Yes      & By clip  & Shaw et al.\ (2018); Transformer-XL, T5, DeBERTa use related but distinct relative schemes \\
RoPE               & None    & Yes      & Good (w/ scaling) & LLaMA, Mistral, Qwen \\
ALiBi              & None    & Yes      & Strong   & BLOOM, MPT \\
\bottomrule
\end{tabularx}
\caption{Comparison of positional encoding methods.  $L$ = max sequence length,
$K$ = max clipping distance, $d$ = model dimension.}
\label{tab:pos-encoding}
\end{table}

\begin{warningbox}
Interviewers frequently ask ``What happens if you give a Transformer no
positional encoding at all?''  The answer: for encoder-only models, the output
is permutation-equivariant and loses all word-order information.  For
causal (decoder-only) models, the causal mask provides \emph{partial}
positional information---the model knows the relative ordering but not the
absolute distance.  Some research (NoPE) has shown surprisingly reasonable
performance with causal masking alone, but this is not standard practice.
\end{warningbox}


% ============================================================================
% SECTION 5: FULL TRANSFORMER ARCHITECTURE
% ============================================================================
\section{The Full Transformer Architecture}
\label{sec:full-arch}

\subsection{Layer Normalization}

Transformers use \textbf{layer normalization} (Ba et al., 2016) rather than
batch normalization, because sequence lengths vary and batch statistics are
unreliable for autoregressive generation.

\[
    \text{LayerNorm}(\vx) = \frac{\vx - \mu}{\sqrt{\sigma^2 + \epsilon}} \odot \bm{\gamma} + \bm{\beta}
\]
where $\mu$ and $\sigma^2$ are computed across the feature dimension for a
single position.

\paragraph{Pre-norm vs.\ Post-norm.}
\begin{itemize}
    \item \textbf{Post-norm (original):} $\text{LN}(\vx + \text{Sublayer}(\vx))$
          ---apply the sublayer first, add the residual, then normalize.
    \item \textbf{Pre-norm (modern standard):}
          $\vx + \text{Sublayer}(\text{LN}(\vx))$---normalize first, apply
          the sublayer, then add the residual.
\end{itemize}

Pre-norm is now the dominant choice because it produces more stable gradients
at initialization (the residual stream passes through unchanged), greatly
reduces sensitivity to learning rate warmup (though production LLMs still
use warmup), and enables training much deeper transformers.

\paragraph{RMSNorm.}
Many modern LLMs (LLaMA, Mistral, Gemma) replace LayerNorm with
\textbf{RMSNorm}, which omits the mean subtraction:
\[
    \text{RMSNorm}(\vx) = \frac{\vx}{\text{RMS}(\vx)} \odot \bm{\gamma},
    \qquad
    \text{RMS}(\vx) = \sqrt{\frac{1}{d}\sum_{i=1}^{d} x_i^2}
\]
RMSNorm is simpler, slightly faster, and performs comparably to full LayerNorm.

\subsection{Residual Connections}

Every sublayer (attention and FFN) uses a residual (skip) connection:
$\vx + \text{Sublayer}(\vx)$.  Residual connections serve two purposes:
\begin{enumerate}
    \item \textbf{Gradient flow:} Provide a direct gradient path from the loss
          to every layer, preventing vanishing gradients in deep networks.
    \item \textbf{Identity initialization:} At initialization, when sublayer
          outputs are near-zero, the network approximates the identity function,
          making optimization easier.
\end{enumerate}

\subsection{Encoder Block}

A single encoder layer consists of:

\begin{mathresult}{Transformer Encoder Block (Pre-Norm)}
\begin{align*}
    \vh' &= \vx + \text{MultiHead}\bigl(\text{LN}(\vx),\; \text{LN}(\vx),\; \text{LN}(\vx)\bigr) \\[3pt]
    \vh'' &= \vh' + \text{FFN}\bigl(\text{LN}(\vh')\bigr)
\end{align*}
Both self-attention inputs (Q, K, V) come from the \emph{same} source.
\end{mathresult}

\subsection{Decoder Block}

A decoder layer adds two features: \textbf{causal masking} in self-attention
(preventing positions from attending to future tokens) and \textbf{cross-attention}
over the encoder output.

\begin{mathresult}{Transformer Decoder Block (Pre-Norm)}
\begin{align*}
    \vh^{(1)} &= \vy + \text{MaskedMultiHead}\bigl(\text{LN}(\vy),\; \text{LN}(\vy),\; \text{LN}(\vy)\bigr) \\[3pt]
    \vh^{(2)} &= \vh^{(1)} + \text{MultiHead}\bigl(\text{LN}(\vh^{(1)}),\; \mH_{\text{enc}},\; \mH_{\text{enc}}\bigr) \\[3pt]
    \vh^{(3)} &= \vh^{(2)} + \text{FFN}\bigl(\text{LN}(\vh^{(2)})\bigr)
\end{align*}
In cross-attention, queries come from the decoder while keys and values come
from the encoder output $\mH_{\text{enc}}$.  Note that only the
\emph{decoder} stream is normalized inside the block: in standard pre-norm
encoder--decoders (e.g., T5), $\mH_{\text{enc}}$ is the encoder output after
a single final LayerNorm at the top of the encoder stack, not re-normalized
inside every decoder layer.
\end{mathresult}

The \textbf{causal mask} is a lower-triangular matrix applied to the attention
scores (setting future positions to $-\infty$ before $\softmax$), ensuring that
each position can only attend to itself and earlier positions.

\subsection{TikZ: Full Transformer Encoder-Decoder}

\begin{figure}[H]
\centering
\begin{tikzpicture}[
    node distance=0.55cm,
    block/.style={rectangle, draw=deepblue, fill=lightblue,
                  minimum width=3.0cm, minimum height=0.55cm,
                  font=\scriptsize\sffamily, rounded corners=2pt, thick},
    normblock/.style={rectangle, draw=gray!60, fill=gray!8,
                      minimum width=3.0cm, minimum height=0.4cm,
                      font=\scriptsize\sffamily, rounded corners=1pt},
    arr/.style={-{Stealth[length=4pt]}, thick, deepblue!70},
    brace/.style={decorate, decoration={brace, amplitude=6pt, raise=4pt},
                  thick, deepblue!60},
    lbl/.style={font=\tiny\sffamily, text=deepblue!70}
]

% === ENCODER (left) ===
\node[block] (einp) {Input Embedding + PE};

\node[normblock, above=0.45cm of einp] (eln1) {Layer Norm};
\node[block, above=0.35cm of eln1] (emha) {Multi-Head Self-Attention};
\node[lbl, above right=-0.12cm and 0.05cm of emha] {Add \& Norm};

\node[normblock, above=0.45cm of emha] (eln2) {Layer Norm};
\node[block, above=0.35cm of eln2] (effn) {Feed-Forward Network};
\node[lbl, above right=-0.12cm and 0.05cm of effn] {Add \& Norm};

\node[block, above=0.6cm of effn, fill=deepblue!15] (eout) {Encoder Output};

% Encoder arrows
\draw[arr] (einp) -- (eln1);
\draw[arr] (eln1) -- (emha);
\draw[arr] (emha) -- (eln2);
\draw[arr] (eln2) -- (effn);
\draw[arr] (effn) -- (eout);

% Residual connections (encoder)
\draw[arr, gray!50, dashed] (einp.west) -- ++(-0.35,0) |- (emha.west);
\draw[arr, gray!50, dashed] (emha.east) -- ++(0.35,0) |- (effn.east);

% Encoder brace
\draw[brace] ([xshift=-0.6cm]eln1.south west) -- ([xshift=-0.6cm]effn.north west)
    node[midway, left=0.45cm, rotate=90, font=\scriptsize\sffamily\bfseries,
         text=deepblue!60] {$\times N$};

% === DECODER (right) ===
\node[block, right=3.5cm of einp] (dinp) {Output Embedding + PE};

\node[normblock, above=0.45cm of dinp] (dln1) {Layer Norm};
\node[block, above=0.35cm of dln1] (dmha) {Masked Self-Attention};

\node[normblock, above=0.45cm of dmha] (dln2) {Layer Norm};
\node[block, above=0.35cm of dln2] (dca) {Cross-Attention};

\node[normblock, above=0.45cm of dca] (dln3) {Layer Norm};
\node[block, above=0.35cm of dln3] (dffn) {Feed-Forward Network};

\node[block, above=0.6cm of dffn, fill=deepblue!15] (dout) {Linear + Softmax};

% Decoder arrows
\draw[arr] (dinp) -- (dln1);
\draw[arr] (dln1) -- (dmha);
\draw[arr] (dmha) -- (dln2);
\draw[arr] (dln2) -- (dca);
\draw[arr] (dca) -- (dln3);
\draw[arr] (dln3) -- (dffn);
\draw[arr] (dffn) -- (dout);

% Cross-attention from encoder
\draw[arr, deepgreen!70, thick] (eout.east) -- ++(0.6,0) |- (dca.west);
\node[lbl, text=deepgreen!70] at ([xshift=0.8cm, yshift=0.3cm]eout.east) {K, V};

% Residual connections (decoder)
\draw[arr, gray!50, dashed] (dinp.west) -- ++(-0.25,0) |- (dmha.west);
\draw[arr, gray!50, dashed] (dmha.east) -- ++(0.25,0) |- (dca.east);
\draw[arr, gray!50, dashed] (dca.west) -- ++(-0.25,0) |- (dffn.west);

% Decoder brace
\draw[brace] ([xshift=3.55cm]dln1.south east) -- ([xshift=3.55cm]dffn.north east)
    node[midway, right=0.45cm, rotate=-90, font=\scriptsize\sffamily\bfseries,
         text=deepblue!60] {$\times N$};

% Labels
\node[font=\small\sffamily\bfseries, text=deepblue, above=0.15cm of eout] {Encoder};
\node[font=\small\sffamily\bfseries, text=deepblue, above=0.15cm of dout] {Decoder};

\end{tikzpicture}
\caption{Full Transformer encoder--decoder architecture (pre-norm variant).
Residual connections (dashed) and layer normalization are applied around each
sublayer.  The encoder stack is repeated $N$ times; the decoder stack similarly
$N$ times, with cross-attention connecting decoder queries to encoder keys and
values.}
\label{fig:transformer-arch}
\end{figure}


% ============================================================================
% SECTION 6: ARCHITECTURE VARIANTS
% ============================================================================
\section{Architecture Variants}
\label{sec:arch-variants}

The original Transformer is an encoder--decoder model, but three distinct
architectural families have emerged, each dominant for different NLP tasks.

\subsection{Encoder-Only (BERT Family)}

\begin{itemize}
    \item \textbf{Attention:} Bidirectional (every token attends to every other
          token, no masking).
    \item \textbf{Pretraining:} Masked Language Modeling (MLM)---predict masked
          tokens from full bidirectional context.
    \item \textbf{Strengths:} Rich, bidirectional representations ideal for
          understanding tasks.
    \item \textbf{Use cases:} Classification, NER, extractive QA, semantic
          similarity, retrieval (bi-encoders).
    \item \textbf{Examples:} BERT, RoBERTa, DeBERTa, ELECTRA, ALBERT.
\end{itemize}

\subsection{Decoder-Only (GPT Family)}

\begin{itemize}
    \item \textbf{Attention:} Causal (lower-triangular mask; each token attends
          only to itself and previous tokens).
    \item \textbf{Pretraining:} Causal Language Modeling (CLM)---predict the
          next token given all preceding tokens.
    \item \textbf{Strengths:} Natural generation, in-context learning emerges
          at scale, simple architecture.
    \item \textbf{Use cases:} Text generation, chat, code generation, reasoning,
          general-purpose via prompting.
    \item \textbf{Examples:} GPT-2/3/4, LLaMA, Mistral, Claude, Gemini.
\end{itemize}

\subsection{Encoder-Decoder (T5 Family)}

\begin{itemize}
    \item \textbf{Attention:} Bidirectional encoder + causal decoder with
          cross-attention.
    \item \textbf{Pretraining:} Span corruption (mask contiguous spans, decode
          them) or other denoising objectives.
    \item \textbf{Strengths:} Explicit separation of input understanding and
          output generation; strong for structured input-to-output tasks.
    \item \textbf{Use cases:} Translation, summarization, structured output,
          sequence-to-sequence tasks.
    \item \textbf{Examples:} T5, Flan-T5, BART, mBART, UL2.
\end{itemize}

\subsection{Comparison Table}

\begin{table}[H]
\centering
\small
\begin{tabularx}{\textwidth}{@{}l c c c X@{}}
\toprule
\textbf{Property} & \textbf{Encoder-Only} & \textbf{Decoder-Only} & \textbf{Enc-Dec} & \textbf{Notes} \\
\midrule
Attention dir.      & Bidirectional & Causal      & Both          & --- \\
Pretraining         & MLM           & CLM         & Span corrupt. & --- \\
In-context learning & Weak          & Strong      & Moderate      & Emerges at scale in CLM \\
Generation          & Poor          & Native      & Good          & Encoder-only needs hacks \\
Understanding       & Excellent     & Good        & Excellent     & Bidirectional helps \\
KV-cache            & N/A           & Efficient   & Partial       & Decoder-only benefits most \\
Scaling trend       & Plateaued     & Dominant    & Niche         & Decoder-only won 2023+ \\
\bottomrule
\end{tabularx}
\caption{Comparison of the three transformer architecture families.}
\label{tab:arch-comparison}
\end{table}

\begin{keyinsight}[Why Decoder-Only Dominated]
The LLM era converged on decoder-only architectures for several reinforcing
reasons: (1)~\textbf{Simplicity}---one stack instead of two reduces
engineering and optimization complexity. (2)~\textbf{Autoregressive scaling}---CLM
naturally benefits from more data (the entire internet is a training corpus).
(3)~\textbf{In-context learning}---the ability to perform tasks from examples in
the prompt emerged specifically in large autoregressive models.
(4)~\textbf{KV-cache efficiency}---causal masking means previously computed
keys and values never change, enabling efficient caching for generation.
Encoder--decoder models remain strong for specific structured tasks
(translation, summarization) but the flexibility of decoder-only models proved
decisive at scale.
\end{keyinsight}


% ============================================================================
% SECTION 7: BERT DEEP DIVE
% ============================================================================
\section{BERT: A Deep Dive}
\label{sec:bert}

BERT (Devlin et al., NAACL 2019) was the first model to demonstrate that
\emph{bidirectional} pretraining followed by task-specific fine-tuning could
achieve state-of-the-art results across a wide range of NLP tasks.

\subsection{Masked Language Modeling (MLM)}

BERT randomly selects 15\% of input tokens for prediction.  Of these selected
tokens:

\begin{itemize}
    \item \textbf{80\% replaced with \texttt{[MASK]}}---the standard masking
          signal.
    \item \textbf{10\% replaced with a random token}---forces the model to
          maintain distributional representations for all positions, not just
          masked ones.  If only \texttt{[MASK]} were used, the model would
          learn that unmasked positions need not be represented well.
    \item \textbf{10\% kept unchanged}---biases the model toward the actual
          observed distribution.  Without this, the model would learn that any
          unmodified token is definitely correct, creating a training-inference
          mismatch.
\end{itemize}

\begin{warningbox}
A common interview mistake is saying ``15\% of tokens are replaced with
\texttt{[MASK]}.''  Only 80\% of the selected 15\% become \texttt{[MASK]}
(i.e., 12\% of all tokens).  Knowing the exact breakdown and \emph{why each
fraction exists} signals depth.
\end{warningbox}

\subsection{Next Sentence Prediction (NSP)}

BERT was also trained to predict whether sentence B follows sentence A (50\%
real next sentences, 50\% random).  The \texttt{[CLS]} token's representation
was used for this binary classification.

RoBERTa (Liu et al., 2019) showed that NSP \emph{hurts} performance by
contaminating the MLM signal, and removed it.  The likely reason: NSP is too
easy---topic detection suffices---and the model wastes capacity on this trivial
signal.

\subsection{Fine-Tuning BERT}

\begin{itemize}
    \item \textbf{Classification:} Add linear head on \texttt{[CLS]} output.
    \item \textbf{NER / Token classification:} Add linear head per token output.
    \item \textbf{Extractive QA:} Predict start and end token positions from
          the token-level outputs over the passage.
    \item \textbf{Sentence similarity:} Siamese BERT (Sentence-BERT) with
          contrastive or cosine loss on \texttt{[CLS]} pooling.
\end{itemize}

\subsection{BERT Variants and Improvements}

\begin{table}[H]
\centering
\small
\begin{tabularx}{\textwidth}{@{}l l X@{}}
\toprule
\textbf{Model} & \textbf{Key Change} & \textbf{Why It Matters} \\
\midrule
RoBERTa     & No NSP, dynamic masking, more data, larger batches
            & Showed BERT was severely undertrained; became the practical default \\
DeBERTa     & Disentangled attention (separate content \& position)
            & SOTA on SuperGLUE; better fine-grained position-content interaction \\
ELECTRA     & Replaced Token Detection (RTD) instead of MLM
            & All tokens provide training signal (not just 15\%); far more efficient \\
SpanBERT    & Mask contiguous spans + span boundary objective
            & Better for span-based tasks (QA, NER, coreference) \\
ALBERT      & Factorized embedding + cross-layer parameter sharing
            & 18$\times$ fewer params with minimal quality drop \\
DistilBERT  & Knowledge distillation from BERT-base
            & 60\% of params, 97\% of performance, 60\% faster \\
\bottomrule
\end{tabularx}
\caption{Key BERT variants and their contributions.}
\label{tab:bert-variants}
\end{table}

\paragraph{ELECTRA in detail.}  Instead of masking tokens, ELECTRA uses a
small generator to \emph{replace} tokens, then trains the main model as a
\textbf{discriminator} to detect which tokens were replaced.  Since \emph{every}
token receives a training signal (binary: real or replaced), ELECTRA is
dramatically more sample-efficient than MLM, which only trains on the 15\%
that are masked.


% ============================================================================
% SECTION 8: GPT FAMILY
% ============================================================================
\section{The GPT Family}
\label{sec:gpt}

The GPT (Generative Pre-trained Transformer) line from OpenAI traces the
evolution from ``pretraining helps'' to ``scale is all you need'' to
``alignment makes models useful.''

\subsection{The Evolution}

\begin{table}[H]
\centering
\small
\begin{tabularx}{\textwidth}{@{}l c l X@{}}
\toprule
\textbf{Model} & \textbf{Year} & \textbf{Params} & \textbf{Key Insight} \\
\midrule
GPT-1    & 2018 & 117M  & Unsupervised pretraining + supervised fine-tuning works for NLP \\
GPT-2    & 2019 & 1.5B  & Scale improves zero-shot; language models are multi-task learners \\
GPT-3    & 2020 & 175B  & In-context learning emerges: few-shot via prompt, no gradient updates \\
InstructGPT & 2022 & 175B & SFT + RLHF: alignment makes models follow instructions safely \\
GPT-4    & 2023 & MoE$^*$ & Multimodal (text + vision), dramatically better reasoning \\
\bottomrule
\end{tabularx}
\caption{GPT family evolution.  $^*$Architecture not officially confirmed but
widely believed to be Mixture of Experts.}
\label{tab:gpt-family}
\end{table}

\subsection{Key Architectural Decisions}

All GPT models use:
\begin{itemize}
    \item \textbf{Decoder-only} architecture with causal masking.
    \item \textbf{Pre-norm} layer normalization (starting from GPT-2).
    \item \textbf{Byte-level BPE} tokenization (GPT-2 onward), ensuring no
          out-of-vocabulary tokens for any language or script.
    \item \textbf{Autoregressive training}: minimize
          $\loss = -\sum_{t=1}^{T} \log P(x_t \mid x_1, \ldots, x_{t-1})$.
\end{itemize}

\subsection{Why Autoregressive Pretraining Enables Few-Shot}

In-context learning (ICL) is the ability to perform a task from examples in
the prompt without any parameter updates.  Why does it emerge?

\begin{enumerate}
    \item \textbf{Implicit task learning during pretraining:} The internet
          contains countless instances of implicit input-output patterns
          (Q\&A forums, translation pairs, code comments).  The model learns
          to recognize and continue these patterns.
    \item \textbf{Induction heads (mechanistic explanation):} Olsson et al.\
          (2022) identified specific attention head circuits that perform
          pattern-matching---they find a previous occurrence of a pattern and
          copy the token that followed it.  This two-head circuit (prefix head
          + induction head) is a mechanism for ICL.
    \item \textbf{Scale is necessary:} ICL emerges reliably only beyond a
          certain scale (roughly $>$1B parameters), suggesting it requires
          sufficient capacity to encode the meta-learning circuits.
\end{enumerate}


% ============================================================================
% SECTION 9: ATTENTION EFFICIENCY
% ============================================================================
\section{Attention Efficiency and Modern Variants}
\label{sec:attention-efficiency}

The $O(n^2)$ cost of self-attention becomes prohibitive for long sequences.
Modern systems address this through several orthogonal strategies.

\subsection{Multi-Query Attention (MQA)}

Shazeer (2019) observed that the KV-cache---not compute---is the primary
bottleneck during autoregressive inference.  MQA uses a \textbf{single shared}
key and value head across all query heads:
\[
    \text{head}_i = \text{Attention}(\mQ \mW_Q^{(i)},\; \mK \mW_K,\; \mV \mW_V)
\]
This reduces KV-cache size by a factor of $h$ (number of heads), dramatically
improving inference throughput for memory-bound generation.  PaLM and
Falcon-7B use MQA (Falcon-40B instead uses grouped attention with 8 KV
heads, a GQA-style configuration).

\subsection{Grouped-Query Attention (GQA)}

Ainslie et al.\ (2023) proposed a middle ground: instead of one KV head
(MQA) or $h$ KV heads (standard MHA), use $g$ groups of KV heads, where each
group is shared by $h/g$ query heads.  GQA with $g = 1$ reduces to MQA;
$g = h$ is standard MHA.

\textbf{Why GQA won:} It preserves most of MQA's memory savings ($h/g$ fold
reduction in KV-cache) while recovering almost all of MHA's quality.  LLaMA~2
70B and Mistral use GQA with $g = 8$.

\begin{table}[H]
\centering
\small
\begin{tabularx}{\textwidth}{@{}l c c c X@{}}
\toprule
\textbf{Method} & \textbf{Query Heads} & \textbf{KV Heads} & \textbf{KV-Cache} & \textbf{Used In} \\
\midrule
MHA (standard)   & $h$  & $h$  & $1\times$                & BERT, GPT-3, GPT-4 \\
MQA              & $h$  & $1$  & $1/h$                    & PaLM, Falcon-7B \\
GQA              & $h$  & $g$  & $g/h$                    & LLaMA 2 70B, Mistral \\
\bottomrule
\end{tabularx}
\caption{Comparison of attention variants.  KV-cache size relative to standard MHA.}
\label{tab:attention-variants}
\end{table}

\subsection{Flash Attention}

Flash Attention (Dao et al., NeurIPS 2022) is not an \emph{approximation}---it
computes \textbf{exact standard attention} but restructures the computation to
minimize memory reads/writes to GPU high-bandwidth memory (HBM).

\paragraph{The core insight.}
Standard attention materializes the full $n \times n$ attention matrix in HBM,
then reads it back to compute the weighted sum.  These HBM reads are the
bottleneck---not the arithmetic.  Flash Attention uses \textbf{tiling}: it
processes the attention matrix in blocks that fit in fast SRAM, computing
partial softmax results and accumulating them without ever materializing the
full matrix.

\paragraph{Key results:}
\begin{itemize}
    \item \textbf{Memory:} $O(n)$ instead of $O(n^2)$---the full attention
          matrix is never stored.
    \item \textbf{Speed:} 2--4$\times$ faster training due to reduced HBM
          I/O.
    \item \textbf{Quality:} Identical to standard attention (no approximation).
\end{itemize}

Flash Attention is now the de facto standard in nearly all transformer
training and inference frameworks.

\begin{keyinsight}[Flash Attention Is About I/O, Not Math]
Flash Attention does not change the attention \emph{computation}---it changes
where the computation \emph{happens}.  By keeping data in fast SRAM and never
writing the full attention matrix to slow HBM, it achieves wall-clock speedups
despite performing the same number of floating-point operations.  This is a
systems optimization, not an algorithmic approximation, which is why it
produces \emph{exact} results.
\end{keyinsight}

\subsection{Sliding Window Attention}

Mistral (Jiang et al., 2023) combines Flash Attention with a \textbf{sliding
window} of size $w$: each token attends only to the $w$ most recent tokens
within the same layer.  However, because information propagates through layers,
a token at layer $l$ has an effective receptive field of $l \times w$ tokens,
meaning the full context is accessible through multi-layer stacking.

This reduces per-layer complexity to $O(n \cdot w)$ and KV-cache to
$O(w \cdot d)$ per layer, enabling efficient processing of long sequences.

\subsection{Summary: Efficiency Techniques}

\begin{table}[H]
\centering
\small
\begin{tabularx}{\textwidth}{@{}l l c c X@{}}
\toprule
\textbf{Technique} & \textbf{Mechanism} & \textbf{Exact?} & \textbf{Speedup} & \textbf{Used In} \\
\midrule
Flash Attention  & IO-aware tiling in SRAM  & Yes & 2--4$\times$ & Nearly all modern training \\
MQA              & Shared K/V              & Yes & 1.5--2$\times$ (inference) & PaLM, Falcon-7B \\
GQA              & Grouped K/V             & Yes & 1.3--1.8$\times$ (inference) & LLaMA 2, Mistral \\
Sliding Window   & Local attention window   & No$^\dagger$ & $O(nw)$ vs.\ $O(n^2)$ & Mistral \\
Sparse Attention & Structured sparsity     & No  & 2--5$\times$ & BigBird, Longformer \\
Linear Attention & Kernel approximation    & No  & $O(n)$ & Research models \\
\bottomrule
\end{tabularx}
\caption{Summary of attention efficiency techniques.  $^\dagger$Sliding window is
exact within the window but does not compute global attention in a single layer.}
\label{tab:efficiency-summary}
\end{table}


% ============================================================================
% SECTION 10: INTERVIEW QUESTIONS
% ============================================================================
\section{Interview Questions}
\label{sec:transformer-questions}

%% QUESTION 1 %%
\begin{interviewq}
\textbf{Q: Explain self-attention step by step, including the math.}

\badgeLfive{L5} \quad \badgetype{Mathematical} \quad \badgefreq{\freqhigh}

\stronganswer
Given input $\mX \in \R^{n \times d}$, we compute three projections:
$\mQ = \mX\mW_Q$, $\mK = \mX\mW_K$, $\mV = \mX\mW_V$.  The attention output
is $\softmax(\mQ\mK^\top / \sqrt{d_k})\,\mV$.  The QK dot product computes
pairwise compatibility scores.  We scale by $\sqrt{d_k}$ because the dot
product variance grows with $d_k$, which would saturate the softmax.  After
softmax normalization (row-wise), each position gets a weighted combination of
all value vectors, where the weights reflect relevance.  Separate Q/K/V
projections allow attending based on one aspect while retrieving another.
Complexity is $O(n^2 d)$ compute and $O(n^2)$ memory for the attention matrix.

\redflags
\begin{itemize}
    \item Cannot write the attention formula
    \item Confuses self-attention with cross-attention
    \item Forgets to mention the scaling factor or cannot explain why it exists
    \item Says ``attention is just a weighted average'' without explaining how
          weights are computed
\end{itemize}

\interviewtest{Whether the candidate can derive the core transformer mechanism
from first principles, not just recite it.  Can they explain \emph{why} each
component (Q/K/V, scaling, softmax) is there?}

\goodgreat
\begin{itemize}
    \item \badgeLfive{L5} States the formula correctly, explains Q/K/V roles,
          mentions complexity
    \item \badgeLsix{L6} Derives the $\sqrt{d_k}$ scaling from the variance
          argument, discusses the attention-as-soft-lookup analogy, connects to
          KV-cache
    \item \badgeLseven{L7} Discusses attention as kernel smoothing, connects to
          Hopfield networks or associative memory, analyzes what information
          different layers capture
\end{itemize}

\followups{How does causal masking change the computation?  What is the
relationship between attention and convolution?  How do you compute attention
gradients efficiently?}
\end{interviewq}

\bigskip

%% QUESTION 2 %%
\begin{interviewq}
\textbf{Q: Why do we scale by $\sqrt{d_k}$? Derive the variance argument.}

\badgeLfive{L5} \quad \badgetype{Mathematical} \quad \badgefreq{\freqhigh}

\stronganswer
Assume $q_i, k_i \sim \mathcal{N}(0, 1)$ i.i.d.  The dot product
$\vq^\top\vk = \sum_{i=1}^{d_k} q_i k_i$ has mean 0 and variance
$\Var(\vq^\top\vk) = \sum_i \Var(q_i k_i) = \sum_i \E[q_i^2]\E[k_i^2] = d_k$.
So the standard deviation is $\sqrt{d_k}$.  For $d_k = 64$, raw dot products
have std $\approx 8$.  The softmax of values in range $[-24, 24]$ is
essentially one-hot, producing vanishing gradients.  Dividing by $\sqrt{d_k}$
normalizes the variance back to 1, keeping softmax in its sensitive region.

\redflags
\begin{itemize}
    \item ``It normalizes the values'' without quantitative reasoning
    \item Cannot derive that the variance equals $d_k$
    \item Confuses this with layer normalization
\end{itemize}

\interviewtest{Mathematical rigor and ability to reason about distributions.
Can the candidate do a simple variance calculation under pressure?}

\goodgreat
\begin{itemize}
    \item \badgeLfive{L5} Correctly derives variance = $d_k$, explains softmax
          saturation
    \item \badgeLsix{L6} Discusses the connection to temperature in softmax,
          notes that this is equivalent to attention temperature $\tau = \sqrt{d_k}$
    \item \badgeLseven{L7} Connects to the entropy of the attention distribution
          and discusses how scaling affects the effective number of tokens attended to
\end{itemize}

\followups{What happens if you use a different scaling factor?  What is the
relationship between scaling and softmax temperature?}
\end{interviewq}

\bigskip

%% QUESTION 3 %%
\begin{interviewq}
\textbf{Q: Why do we need multiple attention heads? What happens with just one head?}

\badgeLfive{L5} \quad \badgetype{Conceptual} \quad \badgefreq{\freqhigh}

\stronganswer
A single head can capture only one attention pattern per layer.  Multiple heads
allow parallel attention in different representation subspaces: one head might
track syntactic dependencies (subject-verb), another positional adjacency,
another semantic similarity.  Each head projects into $d_k = d/h$ dimensions
with independent $\mW_Q^{(i)}, \mW_K^{(i)}, \mW_V^{(i)}$.  The output
projection $\mW_O$ mixes information across heads.  Computationally, multi-head
costs the same as single-head at full dimension because the total dimension is
split across heads.  Empirically, single-head attention at full dimension
performs significantly worse, especially on tasks requiring diverse reasoning.

\redflags
\begin{itemize}
    \item Cannot explain what ``different subspaces'' means concretely
    \item Thinks multi-head attention is more expensive than single-head
    \item Does not mention the output projection $\mW_O$ and its role
\end{itemize}

\interviewtest{Understanding of how multi-head attention provides model capacity
without additional cost, and whether the candidate can reason about what
different heads learn.}

\goodgreat
\begin{itemize}
    \item \badgeLfive{L5} Explains subspace decomposition and the diversity argument
    \item \badgeLsix{L6} References empirical analysis of head behaviors (Voita et al.,
          2019), discusses head pruning experiments that show many heads are redundant
    \item \badgeLseven{L7} Discusses the relationship to mixture-of-experts, the
          lottery ticket hypothesis applied to attention heads, and recent work on
          head importance metrics
\end{itemize}

\followups{How would you determine which heads are important? Can you prune
attention heads?  What is the effect of head dimension on performance?}
\end{interviewq}

\bigskip

%% QUESTION 4 %%
\begin{interviewq}
\textbf{Q: Compare encoder-only, decoder-only, and encoder-decoder transformers.
Why did decoder-only dominate the LLM era?}

\badgeLfive{L5} \quad \badgetype{Conceptual} \quad \badgefreq{\freqhigh}

\stronganswer
\textbf{Encoder-only} (BERT): bidirectional attention, trained with MLM, excels
at understanding/classification tasks but cannot generate text natively.
\textbf{Decoder-only} (GPT): causal attention, trained autoregressively,
natural for generation and showed emergent in-context learning at scale.
\textbf{Encoder-decoder} (T5): bidirectional encoder feeds into causal decoder
via cross-attention, strong for structured input-output tasks like translation.

Decoder-only dominated because: (1)~architectural simplicity---one stack, not
two; (2)~CLM scales naturally with internet-scale data; (3)~in-context learning
emerged, eliminating the need for task-specific fine-tuning; (4)~KV-cache
works cleanly for causal models---previously computed KVs never change.
Encoder-decoder still wins for some structured tasks but decoder-only's
flexibility and scaling properties proved decisive.

\redflags
\begin{itemize}
    \item Cannot articulate the attention masking difference
    \item Says ``GPT is better than BERT'' without nuance
    \item Does not mention in-context learning as a key factor
    \item Confuses training objectives (MLM vs.\ CLM)
\end{itemize}

\interviewtest{Depth of understanding of architectural tradeoffs and ability to
reason about why the field evolved in a particular direction---not just what
happened but why.}

\goodgreat
\begin{itemize}
    \item \badgeLfive{L5} Correctly describes all three, gives examples and use cases
    \item \badgeLsix{L6} Discusses how the training objective (MLM vs.\ CLM) shapes
          capabilities, explains the KV-cache argument in detail
    \item \badgeLseven{L7} Considers hybrid approaches (PrefixLM, UL2), discusses
          whether encoder-decoder could return with the right training recipe, and
          connects to efficiency (FLOPs-per-token vs.\ FLOPs-per-parameter)
\end{itemize}

\followups{When would you still choose BERT over an LLM?  What is prefix
language modeling?  Could encoder-decoder make a comeback?}
\end{interviewq}

\bigskip

%% QUESTION 5 %%
\begin{interviewq}
\textbf{Q: What is the computational complexity of self-attention? Why is it a problem?}

\badgeLfive{L5} \quad \badgetype{Mathematical} \quad \badgefreq{\freqhigh}

\stronganswer
The $\mQ\mK^\top$ product is $\R^{n \times d_k} \cdot \R^{d_k \times n} = \R^{n \times n}$,
requiring $O(n^2 d_k)$ FLOPs.  The subsequent multiplication with $\mV$ is also
$O(n^2 d_v)$.  Total: $O(n^2 d)$ compute and $O(n^2 + nd)$ memory.  For a
128K-token context with 32 heads, the attention matrix alone is
$32 \times 128\text{K}^2 \times 2$ bytes $\approx$ 1\,TB in float16---clearly
impractical without optimizations like Flash Attention (which reduces memory to
$O(n)$ without approximation) or architectural changes like sliding window
attention ($O(nw)$) or linear attention ($O(nd)$).  The quadratic scaling means
doubling context length quadruples both compute and memory, fundamentally
limiting the sequences transformers can process.

\redflags
\begin{itemize}
    \item Says ``linear complexity'' or cannot state $O(n^2)$
    \item Confuses compute complexity with memory complexity
    \item Does not connect to practical implications (context length limits)
\end{itemize}

\interviewtest{Whether the candidate understands the fundamental bottleneck of
transformers and can reason quantitatively about computational costs.}

\goodgreat
\begin{itemize}
    \item \badgeLfive{L5} States $O(n^2 d)$ compute and $O(n^2)$ memory, explains
          why this limits context length
    \item \badgeLsix{L6} Distinguishes between training (Flash Attention helps) and
          inference (KV-cache, MQA/GQA help), discusses the I/O vs.\ compute tradeoff
    \item \badgeLseven{L7} Discusses sub-quadratic alternatives (Mamba, linear attention),
          their tradeoffs vs.\ full attention, and when approximations are acceptable
\end{itemize}

\followups{How does Flash Attention achieve $O(n)$ memory?  What are the
tradeoffs of linear attention?  How does Mamba compare?}
\end{interviewq}

\bigskip

%% QUESTION 6 %%
\begin{interviewq}
\textbf{Q: Explain RoPE (Rotary Position Embedding). Why is it better than
learned absolute positions for long sequences?}

\badgeLsix{L6} \quad \badgetype{Conceptual} \quad \badgefreq{\freqmed}

\stronganswer
RoPE encodes position by rotating query and key vectors in 2D subspaces: for
position $m$ and dimension pair $(2i, 2i{+}1)$, the vector is rotated by angle
$m\theta_i$ where $\theta_i = 10000^{-2i/d}$.  The critical property is that
the inner product of rotated Q and K depends only on the \emph{relative}
position $m - n$: $\text{RoPE}(\vq, m)^\top \text{RoPE}(\vk, n) = f(\vq, \vk, m{-}n)$.
So RoPE achieves relative position encoding through an absolute mechanism,
requiring no extra parameters.  For long sequences, it beats learned absolute
positions because: (1)~no fixed max length---the rotation is defined for any
position; (2)~natural distance decay---inner products decrease with distance;
(3)~length extrapolation via frequency scaling (NTK-aware, YaRN) lets models
trained on 4K context extend to 32K+ by adjusting the base frequency.

\redflags
\begin{itemize}
    \item Cannot explain the rotation mechanism at any level of detail
    \item Confuses RoPE with sinusoidal encoding
    \item Cannot articulate \emph{why} it enables length extrapolation
\end{itemize}

\interviewtest{Whether the candidate understands modern LLM architectural
choices beyond surface level.  RoPE is used in most production LLMs and
understanding it signals genuine depth.}

\goodgreat
\begin{itemize}
    \item \badgeLfive{L5} Describes RoPE as rotating Q/K vectors, mentions relative
          position property
    \item \badgeLsix{L6} Explains the rotation formula, derives the relative position
          property from the inner product, discusses NTK-aware scaling
    \item \badgeLseven{L7} Discusses the connection between RoPE and complex
          exponentials, position interpolation methods (PI, NTK, YaRN), and the
          frequency spectrum interpretation
\end{itemize}

\followups{How does NTK-aware scaling work?  What is YaRN?  How do you extend a
model trained on 4K context to 128K?}
\end{interviewq}

\bigskip

%% QUESTION 7 %%
\begin{interviewq}
\textbf{Q: What does the FFN layer do in a transformer? Is there evidence it
stores factual knowledge?}

\badgeLsix{L6} \quad \badgetype{Conceptual} \quad \badgefreq{\freqmed}

\stronganswer
The FFN is a position-wise two-layer MLP applied independently to each token.
Unlike attention, which mixes tokens, FFN transforms each token's
representation without cross-position interaction.  It expands to $4d$ dimensions
then compresses back, applying nonlinear transformations.  Geva et al.\ (2021)
provided evidence that FFN layers act as key-value memories: rows of $\mW_1$
serve as keys detecting input patterns, and corresponding columns of $\mW_2$
store associated information retrieved when patterns match.  This is supported
by: (1)~knowledge editing techniques like ROME and MEMIT successfully modify
factual associations by editing specific FFN weights; (2)~larger FFNs
(wider or more layers) correlate with better factual recall; (3)~probing
experiments show factual information localizable to specific FFN neurons.
Attention does the routing; FFN stores and processes knowledge.

\redflags
\begin{itemize}
    \item Dismisses FFN as ``just a nonlinearity''
    \item Cannot distinguish FFN's role from attention's role
    \item Unaware of the key-value memory interpretation or knowledge editing work
\end{itemize}

\interviewtest{Whether the candidate understands transformer internals beyond
the attention mechanism.  Most candidates focus on attention and neglect FFN---
knowing the key-value memory interpretation is a strong signal.}

\goodgreat
\begin{itemize}
    \item \badgeLfive{L5} Describes the expand-compress architecture and says FFN
          applies per position
    \item \badgeLsix{L6} Explains the key-value memory interpretation with Geva
          et al., mentions knowledge editing (ROME/MEMIT)
    \item \badgeLseven{L7} Discusses SwiGLU, the relationship between FFN width
          and knowledge capacity, MoE as a way to increase FFN capacity efficiently,
          and the open question of how knowledge is distributed across layers
\end{itemize}

\followups{How does ROME work?  What is the difference between editing
attention vs.\ FFN weights?  How does MoE relate to FFN capacity?}
\end{interviewq}

\bigskip

%% QUESTION 8 %%
\begin{interviewq}
\textbf{Q: Explain BERT's pretraining objectives. Why was NSP removed in
RoBERTa?}

\badgeLfive{L5} \quad \badgetype{Conceptual} \quad \badgefreq{\freqhigh}

\stronganswer
BERT uses two pretraining objectives: (1)~\textbf{MLM}: 15\% of tokens are
selected; of those, 80\% become \texttt{[MASK]}, 10\% random, 10\% unchanged.
The model predicts original tokens from bidirectional context.  The 80/10/10
split prevents mismatch between pretraining (where \texttt{[MASK]} appears)
and fine-tuning (where it does not).  (2)~\textbf{NSP}: binary prediction of
whether two sentences are consecutive.  RoBERTa showed NSP hurts by removing
it and achieving consistently better results.  The likely reasons: NSP is a
trivially easy task (topic matching suffices), it wastes model capacity on a
useless signal, and the ``random sentence'' negatives are too easy (different
documents = different topics).  RoBERTa also found that dynamic masking
(re-masking each epoch), larger batch sizes, and much more training data
(160GB vs.\ 16GB) significantly improved results---BERT was simply
undertrained.

\redflags
\begin{itemize}
    \item Says ``15\% of tokens become \texttt{[MASK]}'' (only 12\% actually do)
    \item Cannot explain the rationale for the 80/10/10 split
    \item Does not know that RoBERTa removed NSP or cannot say why
\end{itemize}

\interviewtest{Detailed understanding of a foundational model.  The 80/10/10
rationale and NSP critique are shibboleths for genuine NLP depth.}

\goodgreat
\begin{itemize}
    \item \badgeLfive{L5} Correctly explains both objectives with the 80/10/10
          split and NSP removal rationale
    \item \badgeLsix{L6} Discusses alternative objectives (ELECTRA's RTD,
          SpanBERT's span masking) and why they improve on MLM
    \item \badgeLseven{L7} Compares MLM and CLM as pretraining objectives from
          an information-theoretic perspective, discusses bidirectional vs.\
          unidirectional context and the implications for downstream tasks
\end{itemize}

\followups{What is ELECTRA and why is it more efficient?  Why does masking
ratio of 15\% work?  What happens if you mask 50\%?}
\end{interviewq}

\bigskip

%% QUESTION 9 %%
\begin{interviewq}
\textbf{Q: What is Flash Attention? Why does it help if it computes the same
result as standard attention?}

\badgeLsix{L6} \quad \badgetype{Systems} \quad \badgefreq{\freqmed}

\stronganswer
Flash Attention computes \emph{exact} standard attention but restructures the
computation to minimize HBM (high-bandwidth memory) I/O.  Standard attention
materializes the full $n \times n$ attention matrix in HBM, requiring $O(n^2)$
reads and writes.  Modern GPUs are memory-bandwidth-bound, not compute-bound,
for attention.  Flash Attention uses \textbf{tiling}: it processes Q, K, V in
blocks that fit in fast on-chip SRAM, computes partial softmax results using
the log-sum-exp trick for numerical stability, and accumulates without ever
writing the full matrix to HBM.  Result: $O(n)$ memory (no materialized
attention matrix), 2--4$\times$ wall-clock speedup, and identical output.
Flash Attention 2 further improved by reducing non-matmul FLOPs and better
work partitioning.  It is now the default in virtually all transformer
training and serving frameworks.

\redflags
\begin{itemize}
    \item Thinks Flash Attention is an approximation (it is exact)
    \item Cannot explain the HBM vs.\ SRAM distinction
    \item Does not know what ``IO-aware'' means in this context
\end{itemize}

\interviewtest{Understanding of the hardware-software co-design that makes
modern LLMs practical.  This separates candidates who only think at the
algorithm level from those who understand systems.}

\goodgreat
\begin{itemize}
    \item \badgeLfive{L5} Knows Flash Attention is exact and reduces memory; mentions
          tiling
    \item \badgeLsix{L6} Explains the HBM vs.\ SRAM bottleneck, describes the tiling
          strategy, and mentions the online softmax trick
    \item \badgeLseven{L7} Discusses Flash Attention 2/3 improvements, the roofline
          model for analyzing attention performance, and how Flash Attention
          interacts with tensor parallelism across GPUs
\end{itemize}

\followups{How does the online softmax trick work?  What changed in Flash
Attention 2?  How does Flash Attention interact with GQA?}
\end{interviewq}

\bigskip

%% QUESTION 10 %%
\begin{interviewq}
\textbf{Q: Compare Multi-Query Attention, Grouped-Query Attention, and standard
Multi-Head Attention. When would you use each?}

\badgeLsix{L6} \quad \badgetype{Trade-off} \quad \badgefreq{\freqmed}

\stronganswer
\textbf{MHA:} Each head has its own Q, K, V projections ($h$ sets of K/V).
Highest quality, but KV-cache scales as $O(h \cdot n \cdot d_k)$ per layer---
the main bottleneck for inference with long sequences or large batches.
\textbf{MQA:} All $h$ query heads share a single K and V head.  Reduces
KV-cache by $h\times$ (e.g., $32\times$ for a 32-head model), dramatically
improving inference throughput.  Quality drops 1--2\% on some benchmarks.
\textbf{GQA:} A middle ground with $g$ KV head groups.  Each group is shared
by $h/g$ query heads.  With $g = 8$ on a 32-head model, KV-cache is $4\times$
smaller than MHA while recovering most quality.  GQA is the current
best practice.

\textbf{When to use each:} MHA when quality is paramount and inference is
not memory-bound (short sequences, small batches, or training only).  MQA for
maximum throughput in long-context serving.  GQA as the default for production
LLMs (quality close to MHA, memory close to MQA).

\redflags
\begin{itemize}
    \item Cannot explain how GQA interpolates between MHA and MQA
    \item Does not connect to KV-cache memory as the motivation
    \item Thinks these only matter for training (they primarily help inference)
\end{itemize}

\interviewtest{Whether the candidate understands the inference bottleneck
(memory bandwidth, not compute) and can reason about production deployment
tradeoffs for LLMs.}

\goodgreat
\begin{itemize}
    \item \badgeLfive{L5} Correctly describes all three variants and their KV-cache
          implications
    \item \badgeLsix{L6} Quantifies the KV-cache savings, discusses uptraining
          from MHA to GQA, mentions which production models use which variant
    \item \badgeLseven{L7} Discusses Multi-Head Latent Attention (DeepSeek-V2),
          the relationship between attention variants and tensor parallelism
          strategies, and how KV-cache compression interacts with quantization
\end{itemize}

\followups{How would you convert an MHA model to GQA after training?  What is
the KV-cache size for a 70B model with 128K context?  How does KV-cache
quantization work?}
\end{interviewq}

\bigskip

%% QUESTION 11 %%
\begin{interviewq}
\textbf{Q: Explain pre-norm vs.\ post-norm in transformers. Why did pre-norm become
the standard?}

\badgeLfive{L5} \quad \badgetype{Conceptual} \quad \badgefreq{\freqmed}

\stronganswer
\textbf{Post-norm} (original Transformer): $\text{LN}(\vx + \text{Sublayer}(\vx))$---the
residual and sublayer output are added, then normalized.  \textbf{Pre-norm}
(GPT-2 and all modern LLMs): $\vx + \text{Sublayer}(\text{LN}(\vx))$---input
is normalized before the sublayer, then the (unnormalized) sublayer output is
added to the residual.  Pre-norm won because: (1)~the residual stream passes
through without any transformation, providing a clean gradient highway from
loss to any layer; (2)~gradients at initialization are well-behaved (sublayer
outputs are near-zero, so the network starts as approximate identity);
(3)~learning rate warmup becomes less critical---post-norm is notoriously
sensitive to warmup schedule; (4)~enables training much deeper models
(100+ layers) without stability issues.  The tradeoff: some evidence suggests
post-norm can achieve slightly higher final performance if properly trained,
but pre-norm's training stability advantage overwhelmingly dominates in practice.

\redflags
\begin{itemize}
    \item Cannot state the difference in formula
    \item Does not explain \emph{why} pre-norm is more stable
    \item Confuses layer normalization with batch normalization
\end{itemize}

\interviewtest{Understanding of training dynamics and the subtle but critical
architectural choices that make deep transformers trainable.}

\goodgreat
\begin{itemize}
    \item \badgeLfive{L5} Correctly states both formulas and the stability argument
    \item \badgeLsix{L6} Discusses the gradient flow analysis, mentions RMSNorm
          as a further simplification, and knows which models use which variant
    \item \badgeLseven{L7} Discusses DeepNorm, the signal propagation theory of
          pre-norm, and emerging work on normalization-free architectures
\end{itemize}

\followups{What is RMSNorm?  Why does post-norm need learning rate warmup?
What is the gradient flow difference mathematically?}
\end{interviewq}

\bigskip

%% QUESTION 12 %%
\begin{interviewq}
\textbf{Q: You need to build a production NLP system that classifies customer
support tickets into 50 categories. Would you use BERT or GPT-4? Justify your
choice.}

\badgeLfive{L5} \quad \badgetype{Trade-off} \quad \badgefreq{\freqmed}

\stronganswer
For this task, a \textbf{fine-tuned BERT} (or RoBERTa/DeBERTa) is almost
certainly the right choice.  Reasoning: (1)~\textbf{Latency}: a fine-tuned
BERT-base processes a ticket in $\sim$2ms on GPU vs.\ 500ms+ for GPT-4 API.
(2)~\textbf{Cost}: BERT inference is $\sim$1000$\times$ cheaper than GPT-4 API
calls at scale.  (3)~\textbf{Quality}: with 1000+ labeled examples per class,
a fine-tuned encoder model will match or exceed GPT-4 few-shot on a well-defined
classification task.  (4)~\textbf{Control}: you own the model, data stays
private, no API dependency.  (5)~\textbf{Consistency}: fine-tuned classifiers
produce deterministic outputs; LLMs are stochastic and may format outputs
inconsistently.

GPT-4 would be preferred only if: (a)~you have near-zero labeled data
(cold-start), (b)~categories change frequently and retraining is impractical,
or (c)~you need explanations alongside classifications.  The answer is not about
which model is ``better'' in isolation---it is about matching the solution to
the constraints.

\redflags
\begin{itemize}
    \item Reflexively says ``use GPT-4 for everything''
    \item Cannot reason about cost, latency, or deployment constraints
    \item Does not consider the labeled data regime
\end{itemize}

\interviewtest{Production thinking and resistance to hype.  Can the candidate
make pragmatic engineering decisions rather than defaulting to the most powerful
model?}

\goodgreat
\begin{itemize}
    \item \badgeLfive{L5} Recommends BERT with clear cost/latency/quality reasoning
    \item \badgeLsix{L6} Discusses the labeled data regime, considers a hybrid
          approach (GPT-4 for labeling seed data, BERT for serving), and mentions
          distillation from LLM to smaller model
    \item \badgeLseven{L7} Designs a full system: LLM for cold-start and edge
          cases, distilled classifier for the head, active learning loop for
          continuous improvement, monitoring for distribution shift
\end{itemize}

\followups{How would you handle a cold-start with zero labeled data?  How would
you detect when categories shift?  How would you monitor classifier quality in
production?}
\end{interviewq}


% ============================================================================
% CHAPTER SUMMARY
% ============================================================================
\section*{Chapter Summary}

\begin{enumerate}
    \item \textbf{Self-attention} computes $\softmax(\mQ\mK^\top/\sqrt{d_k})\mV$,
          enabling direct interaction between all positions in $O(n^2 d)$ time.
          The scaling factor $\sqrt{d_k}$ prevents softmax saturation.

    \item \textbf{Multi-head attention} runs $h$ parallel attention heads in
          different subspaces, capturing diverse relationships at no extra cost.

    \item \textbf{FFN layers} transform each position independently, acting as
          key-value memories that store factual knowledge.  Modern variants
          (SwiGLU, GeGLU) use gating for improved performance.

    \item \textbf{Positional encoding} is necessary because attention is
          permutation-equivariant.  RoPE is the current standard, encoding
          relative position through rotation of Q/K vectors.

    \item \textbf{The full architecture} combines multi-head attention, FFN,
          residual connections, and layer normalization.  Pre-norm with RMSNorm
          is the modern standard.

    \item \textbf{Three families} emerged: encoder-only (BERT, for understanding),
          decoder-only (GPT, for generation and in-context learning), and
          encoder--decoder (T5, for structured tasks).  Decoder-only dominates
          the LLM era.

    \item \textbf{BERT} introduced bidirectional pretraining with MLM (80/10/10
          masking).  RoBERTa showed it was undertrained and NSP was harmful.
          ELECTRA improved efficiency via replaced token detection.

    \item \textbf{GPT} models showed that scale enables zero-shot and few-shot
          capabilities through in-context learning, and that alignment (RLHF)
          makes models practically useful.

    \item \textbf{Efficiency improvements}---Flash Attention (IO-aware tiling),
          GQA (grouped key-value heads), and sliding window attention---make
          transformers practical for long sequences and production serving.
\end{enumerate}
